\documentclass[12pt,a4paper,oneside]{article}
\usepackage[brazil]{babel}
\usepackage[utf8]{inputenc}
\usepackage{olymp}
\usepackage{color}
\usepackage[dvipsnames]{xcolor}
\usepackage{listings}

\setcounter{problem}{1}

\begin{document}

\begin{problem}{Primos}{primos}{2}
 
O seguinte programa foi escrito para gerar uma lista de números primos. Após ler dois valores inteiros $A$ e $B$, ele deveria gerar todos os números primos entre $A$ e $B$, inclusive.

\lstset{language=C++,
                basicstyle=\ttfamily,
                keywordstyle=\color{Mulberry}\ttfamily,
                stringstyle=\color{Maroon}\ttfamily,
                commentstyle=\color{YellowGreen}\ttfamily,
                morecomment=[l][\color{magenta}]{\#},
                basewidth = {.48em},
                showstringspaces=false
}
{\small
\begin{lstlisting}
#include <iostream>
using namespace std;

int main()
{
    int a, b;

    cin >> a >> b;
    
    cout << "Primos entre " << a << " e " << b << ":";
    for(int i = a; i <= b; i++)
        if (primo(i))           //Se for primo
            cout << " " << i;   //imprime o valor
    cout << endl;

    return 0;
}
\end{lstlisting}
}

O problema é que não existe a função \texttt{primo} usada neste programa. Sua tarefa é completar o programa criando esta função. A função deve seguir o cabeçalho abaixo. Ela recebe um número inteiro e retorna true/false caso o número recebido seja primo ou não\footnote{Se fizer na linguagem C, a função deve ser \texttt{int primo(int n)} e retornar 1 ou 0, pois em C não há \texttt{bool}.}.

\texttt{\textcolor{Mulberry}{bool} primo(\textcolor{Mulberry}{int} n)}

%\Notes
%
%Observações, se tiver.

\InputFile

A entrada contém dois valores inteiros, $A$ e $B$. Restrição: $1 \leq A \leq B \leq 1000$.


\OutputFile

Seu programa deve gerar apenas uma linha de saída, contendo a lista de primos entre $A$ e $B$, escritos conforme os exemplos a seguir.

\Notes
Nesta questão, seu programa deve usar a função \texttt{main} dada acima exatamente como está, não a modifique! Sua tarefa é apenas criar a função \texttt{primo}.

\Example

\begin{example}
\exmp{
1 10
}{
Primos entre 1 e 10: 2 3 5 7
}%
\end{example}

\begin{example}
\exmp{
2 11
}{
Primos entre 2 e 11: 2 3 5 7 11
}%
\end{example}

\begin{example}
\exmp{
315 340
}{
Primos entre 315 e 340: 317 331 337
}%
\end{example}


%Comentários sobre o exemplo, se necessário.

\end{problem}

\end{document}
